\subsection{Convolution divided powers and the two coalgebra questions}

We formulate the naturality implicit in scans 324--335 explicitly.
For a counital cocommutative $k$-coalgebra $(A,\Delta,\epsilon)$ we
seek PD laws on
\[
 \operatorname{Hom}_k(A,J)\subset\operatorname{Hom}_k(A,B),
 \qquad f*g=\mu_B(f\otimes g)\Delta,
\]
for every PD $k$-algebra $(B,J,\gamma)$, natural for PD algebra maps.
No flatness is assumed unless stated.

\begin{theorem}[D007: \status{Proved}; explicit classification of natural laws]
\label{sa08:D007}
Define the following canonical maps:
\begin{align*}
 \delta_{r,s}&:\Gamma^{r+s}A\longrightarrow
                  \Gamma^rA\otimes\Gamma^sA,\\
 \rho_{m,n}&:\Gamma^m(\Gamma^nA)\longrightarrow\Gamma^{mn}A,\\
 \chi_n&:\Gamma^n(A\otimes A)\longrightarrow
                    \operatorname{Sym}^nA\otimes\Gamma^nA,\\
 \pi_n&:A^{\otimes n}\longrightarrow\operatorname{Sym}^nA.
\end{align*}
Here $\delta$ is the homogeneous component of $\Gamma$ of the
additive diagonal, $\rho$ is induced by taking divided powers in
$\Gamma(A)$, and $\chi_n$ is the universal divided power of
$a\otimes b\mapsto a\otimes\gamma_1(b)$ in
$\operatorname{Sym}(A)\otimes\Gamma(A)$. Put
$c_{m,n}=(mn)!/(m!(n!)^m)$. Let $\mathcal P(A)$ consist of the families
$u_n\in\operatorname{Hom}_k(A,\Gamma^nA)$ satisfying
\begin{align}
 u_0&=\epsilon,&u_1&=\mathrm{id},\tag{U}\label{eq:D-U}\\
 \delta_{r,s}u_{r+s}&=(u_r\otimes u_s)\Delta,
       &&r,s\geq0,\tag{A}\label{eq:D-A}\\
 \chi_n\Gamma^n(\Delta)u_n
   &=(\pi_n\otimes u_n)\Delta^{(n+1)},
       &&n\geq1,\tag{S}\label{eq:D-S}\\
 \rho_{m,n}\Gamma^m(u_n)u_m&=c_{m,n}u_{mn},
       &&m,n\geq1.\tag{I}\label{eq:D-I}
\end{align}
Then $\mathcal P(A)$ is exactly the set of natural convolution PD
laws. Its point $u$ gives
\[
 \gamma_n^A(f)=\gamma_{J,n}\,\Gamma^n(f)\,u_n,
 \qquad f:A\to J,
\]
where $\gamma_{J,n}:\Gamma^nJ\to J$ represents the given divided
power. Thus the class of coalgebras requested in B.7 is explicitly
the class with $\mathcal P(A)\ne\varnothing$, and $\mathcal P(A)$
records its entire indeterminacy, including torsion.
\end{theorem}
\begin{proof}
The functor $(B,J)\mapsto\operatorname{Hom}_k(A,J)$ is represented
by the free PD algebra $\Gamma(A)$. A natural operation is therefore
determined by its value on $A\to\Gamma^1A$, a linear map
$A\to\Gamma^+A$. Apply the scalar axiom to $z\epsilon$ after base
extension to $k[z]$. Comparing polynomial coefficients in
$\gamma_n^A(zf)=z^n\gamma_n^A(f)$ forces weight $n$; thus the value
is precisely $u_n:A\to\Gamma^nA$. Evaluation under the universal
PD map to $(B,J)$ gives the displayed formula and proves its
necessity and uniqueness.

Test the sum axiom on the universal pair in $\Gamma(A\oplus A)$.
Under $\Gamma(A\oplus A)=\Gamma(A)\otimes\Gamma(A)$, its
$(r,s)$-component is exactly \eqref{eq:D-A}. Test the scalar axiom
on the universal maps $A\to\operatorname{Sym}(A)$ and
$A\to\Gamma^1A$ in $\operatorname{Sym}(A)\otimes\Gamma(A)$.
This is the free PD algebra over $\operatorname{Sym}(A)$ on
$\operatorname{Sym}(A)\otimes A$, so it is valid without flatness.
Taking the divided power of their convolution gives the left side
of \eqref{eq:D-S}; multiplying the ordinary $n$-fold convolution
power of the first by the new divided power of the second gives
the right side. Universality proves necessity and sufficiency of
this equation.

The product identity follows from \eqref{eq:D-A} by
\[
 \mu\delta_{r,s}=\binom{r+s}{r}\mathrm{id}_{\Gamma^{r+s}A}.
\]
On a divided-power monomial this is Vandermonde's identity, summing
over the choices of $r$ factors; it is integral and hence valid for
all $k$. Finally the iteration axiom, evaluated on the universal
map, is \eqref{eq:D-I}. Equations \eqref{eq:D-U} give the remaining
unit and identity axioms. Every evaluation respects all equations,
so the construction is natural and proves completeness.
\end{proof}

The stronger relation asked on scan 324 is
\begin{equation}\label{eq:D-strong}
 \Gamma^m(u_n)u_m=\sigma_{m,n}u_{mn},
\end{equation}
where $\sigma_{m,n}:\Gamma^{mn}A\to\Gamma^m(\Gamma^nA)$ represents
the polynomial law $a\mapsto\gamma_m(\gamma_n(a))$.
Since $\rho_{m,n}\sigma_{m,n}=c_{m,n}$, it implies the iteration
relation in D007. We keep this condition before $\rho$ distinct
from the full natural PD classification after $\rho$.

\begin{theorem}[D008: \status{Proved}; flat existence and uniqueness]
\label{sa08:D008}
If $A$ is $k$-flat, $\mathcal P(A)$ is a singleton, and its family
also satisfies \eqref{eq:D-strong}. It is
\[
 u_n=\theta_n^{-1}\Delta^{(n)},\qquad
 \theta_n:\Gamma^nA\xrightarrow{\sim}(A^{\otimes n})^{\mathfrak S_n}.
\]
\end{theorem}
\begin{proof}
The flat comparison is Lundkvist (2008), Proposition 2.2, opened
in full HTML and in its arXiv source. A short proof verifies the
form used here: for finite free modules the divided-power basis
maps to the orbit-sum basis of invariant tensors. Every flat
module is a filtered colimit of finite free modules. Divided powers,
tensor powers and invariants of a finite group commute with this
colimit; for the last assertion use exactness of filtered colimits
and the finite equalizer defining invariants. Thus $\theta_n$ is
an isomorphism. The same argument shows that $\Gamma^nA$ is flat.

Coassociativity and cocommutativity make $\Delta^{(n)}$ invariant.
After the injective tensor comparison, equation~\eqref{eq:D-A}
groups an iterated coproduct into two blocks;
\eqref{eq:D-S} selects one factor from each of $n$ pairs; and
\eqref{eq:D-strong} groups $mn$ factors into $m$ blocks of $n$.
Each is exactly coassociativity and cocommutativity. The necessary
iterated comparisons are injective because $\Gamma^nA$ is flat.
Thus this family belongs to $\mathcal P(A)$. Conversely iterating
\eqref{eq:D-A} into $n$ singleton blocks forces
$\theta_nu_n=\Delta^{(n)}$. Injectivity gives uniqueness.
\end{proof}

\begin{proposition}[D009: \status{Refuted}; unrestricted existence]
\label{sa08:D009}
A finite nonflat cocommutative coalgebra can fail to admit any
convolution PD structure, even for one finite PD target.
\end{proposition}
\begin{proof}
Take $k=\mathbb Z/4$ and
\[
 A=k1\oplus(k/2)e,\qquad
 \epsilon(e)=0,\qquad
 \Delta(e)=e\otimes1+1\otimes e+e\otimes e.
\]
The relations respect $2e=0$; counitality and cocommutativity are
immediate. Coassociativity follows because $g=1+e$ is group-like.
The summand $k/2$ is not flat: tensoring $(2)\hookrightarrow k$
with it destroys injectivity.

Let $B=k\oplus(\mathbb Z/2)t\oplus(\mathbb Z/4)u$, with
$t^2=2u$ and $tu=u^2=0$. Its ideal $J=(t,u)$ has
\[
 \gamma_2(at+bu)=a^2u,\qquad\gamma_n=0\quad(n\geq3).
\]
This is well-defined since $(a+2)^2-a^2$ is divisible by $4$.
The sum axiom in degree two is
$(a+c)^2-a^2-c^2=2ac$. For a scalar
$z=\alpha+\beta t+\chi u$,
$z(at+bu)=\alpha at+(\alpha b+2\beta a)u$; hence
$\gamma_2(zx)=\alpha^2a^2u=z^2\gamma_2(x)$.
Also $x^2=2\gamma_2(x)$. Every product involving a member of $J$
and $u$ is zero. Thus all higher sum and product terms vanish,
as do $\gamma_2(\gamma_2(x))$ and all higher iterations. This
verifies every PD axiom, not merely a finite range of indices.

The linear map $f:A\to J$, $f(1)=0$, $f(e)=t$, satisfies
\[
 (f*f)(e)=t^2=2u\ne0.
\]
But any linear $v:A\to B$ has $2v(e)=v(2e)=0$, making
$f*f=2\gamma_2(f)$ impossible. Here
$|A|=8$, $|B|=32$, and $|J|=8$. The supplied program exhaustively
checks the finite ring and PD identities, coalgebra identities and
the obstruction. This is a finite adaptation of the supplied Pro
notes N43, independently verified; no priority is claimed.
\end{proof}

\begin{proposition}[D010: \status{Refuted}; uniqueness from the two written relations]
\label{sa08:D010}
The product relation and \eqref{eq:D-strong} alone do not force the
family $u_n$, even when $A=k=\mathbb F_2$ is flat.
\end{proposition}
\begin{proof}
Besides the canonical family $u_n(1)=\gamma_n(1)$, take
$v_0=1$, $v_1=\mathrm{id}$, and $v_n=0$ for $n\geq2$.
For the product identity the case $r=s=1$ has coefficient $2=0$;
if either index is at least two both sides vanish; zero indices
give the unit. For strong iteration an index one is tautological,
and otherwise both sides vanish. Thus both written relations hold.
But addition compatibility fails at $r=s=1$: the left side of
\eqref{eq:D-A} is zero and the right side is $1\otimes1$.
This family therefore does not give a full natural PD law.
The finite coefficient program checks every pair of indices at
most sixteen; the preceding case argument proves all indices.
The corrected flat uniqueness assertion is D008.
\end{proof}

\begin{theorem}[D011: \status{Proved}; nonflat existence and genuine nonuniqueness]
\label{sa08:D011}
For every module $M$, the primitive coalgebra $A=k\oplus M$ admits
a natural convolution PD law. Full natural PD laws need not be
unique when $M$ is nonflat.
\end{theorem}
\begin{proof}
With $\Delta(m)=m\otimes1+1\otimes m$, convolution is the
square-zero extension $B\ltimes\operatorname{Hom}_k(M,B)$.
Its canonical law is
\[
 \gamma_n(b,x)=(\gamma_n(b),\gamma_{n-1}(b)x).
\]
This is Stacks, Lemma 60.3.1, tag 07HH, opened in this run, applied
to the square-zero extension. Its PD ideal
$J\oplus\operatorname{Hom}(M,B)$ has the stable subideal
$J\oplus\operatorname{Hom}(M,J)$ by the formula, proving existence
without flatness.

For $M=k/2$, write $\operatorname{Hom}(M,B)=B[2]$. A second law,
for $n\geq2$, is
\[
 \gamma'_n(b,x)=
 (\gamma_n(b),\gamma_{n-1}(b)x+\gamma_{n-2}(b)x^2),
 \qquad \gamma'_1=\mathrm{id},\quad\gamma'_0=1.
\]
It is defined for $b\in J$ and $x\in J[2]$. Since $2x=0$,
$4\gamma_2(x)=\gamma_2(2x)=0$, whence
$x^4=(2\gamma_2(x))^2=0$; also $2x^2=0$.
The sum axiom follows by expanding $\gamma_n(b+c)$ and
$(x+y)^2=x^2+y^2$.

For the scalar axiom multiply by $(a,h)\in B\ltimes B[2]$.
The new tangent part is $ax+bh$.
Use $b\gamma_{n-1}(b)=n\gamma_n(b)$ and
$b^2\gamma_{n-2}(b)=n(n-1)\gamma_n(b)$.
The cross term $2abxh$ vanishes, as does the last expression
multiplied by $h^2$. The tangent part of the left side is consequently
\[
 a^n\bigl(\gamma_{n-1}(b)x+\gamma_{n-2}(b)x^2\bigr)
       +na^{n-1}\gamma_n(b)h,
\]
which is precisely that of $(a,h)^n\gamma'_n(b,x)$.
For the product identity, Pascal's identity gives the $x$
coefficient. The difference in the $x^2$ coefficients is
$2\binom{m+n-2}{m-1}$ and annihilates $x^2$.

For iteration with $m,n\geq2$, put
$y=\gamma_{n-1}(b)x+\gamma_{n-2}(b)x^2$.
Then $y^2=\gamma_{n-1}(b)^2x^2$.
The coefficient of $\gamma_{mn-1}(b)x$ is
\[
 \frac{(mn-1)!}{(m-1)!(n!)^{m-1}(n-1)!}=c_{m,n}.
\]
The coefficient of $\gamma_{mn-2}(b)x^2$ is
\begin{align*}
 &\frac{(mn-2)!}{(m-1)!(n!)^{m-1}(n-2)!}\\
 &\hspace{1em}+\frac{(mn-2)!}{(m-2)!(n!)^{m-2}((n-1)!)^2}
 =c_{m,n}.
\end{align*}
Dividing these two integer summands by $c_{m,n}$ gives
$(n-1)/(mn-1)$ and $(m-1)n/(mn-1)$, whose sum is one.
Indices one are tautological. All formulas are natural under PD
algebra maps and agree with the original PD structure on $(b,0)$.

For a distinction take $k=\mathbb Z/4$ and $(B,J)$ from D009,
now with the primitive source coalgebra. At $(0,t)$ the canonical
second divided power is zero and the new one is $(0,2u)\ne0$.
The program exhaustively verifies both structures on the
32-element convolution ideal in its 256-element ring.
Thus flatness is sufficient but not necessary, and full uniqueness
without flatness is false. D007 gives the complete equation
presentation: it is empty in D009, a singleton in D008, and has
at least two points here.
\end{proof}

\begin{theorem}[D012: \status{Proved}; the dual question]
\label{sa08:D012}
Let $B'\to B$ be a map of counital cocommutative coalgebras,
$C=\operatorname{coker}(B'\to B)$ and $\pi:B\to C$.
Divided copowers can be specified without finite duals as maps
$d_n:C\to\operatorname{Sym}^nC$, $n\geq1$, satisfying the
following explicit equations. Put $\widetilde d_n=d_n\pi$
for $n\geq1$ and $\widetilde d_0=\epsilon_B$, and require
$d_1=\mathrm{id}_C$ and
\begin{align}
 \delta_{r,s}\widetilde d_{r+s}
  &=(\widetilde d_r\otimes\widetilde d_s)\Delta_B,
       \tag{DA}\label{eq:D-DA}\\
 \chi_n\operatorname{Sym}^n(\ell)d_n
  &=(\pi_n\otimes d_n)\ell_n,
       \tag{DS}\label{eq:D-DS}\\
 \rho_{m,n}\operatorname{Sym}^m(d_n)d_m
  &=c_{m,n}d_{mn}.\tag{DI}\label{eq:D-DI}
\end{align}
Here $\delta$ is induced by the additive diagonal on $C$;
$\ell:C\to B\otimes C$ and $\ell_n:C\to B^{\otimes n}\otimes C$
are the induced coactions;
$\chi_n:\operatorname{Sym}^n(B\otimes C)\to
\operatorname{Sym}^nB\otimes\operatorname{Sym}^nC$ separates
products of pairs; $\pi_n$ is the symmetric quotient; and $\rho$
multiplies symmetric monomials.
These structures are in bijection with PD laws, natural in all
commutative $k$-algebras $R$, on
\[
 \ker(\operatorname{Hom}_k(B,R)\to\operatorname{Hom}_k(B',R))
 =\operatorname{Hom}_k(C,R).
\]
Their formula is
$\gamma_n(f)=\mu_R^{(n)}\operatorname{Sym}^n(f)d_n$.
No flatness or projectivity hypothesis is required.
\end{theorem}
\begin{proof}
The coaction is well-defined because $(1\otimes\pi)\Delta_B$
vanishes on the image of the coalgebra $B'$; the iterated versions
do also. Evaluating \eqref{eq:D-DA}, \eqref{eq:D-DS} and
\eqref{eq:D-DI} in $R$ gives exactly the sum, scalar and iteration
axioms. The product axiom follows from
$\mu\delta_{r,s}=\binom{r+s}{r}$. Conversely
$\operatorname{Hom}_k(C,R)$ is represented by
$\operatorname{Sym}(C)$. The value of a natural operation on the
universal map determines $d_n$; scaling by $z$ over $k[z]$ forces
symmetric degree $n$. Testing in
$\operatorname{Sym}(C\oplus C)$,
$\operatorname{Sym}(B)\otimes\operatorname{Sym}(C)$, and
$\operatorname{Sym}(C)$ recovers the three displayed equations.
This proves both necessity and sufficiency, including all torsion.

For the example on scan 324, take $B=\operatorname{Sym}_k(M)$
with primitive coproduct and $B'=k$. For $f(1)=0$ define
\[
 \gamma_n(f)(x_1\cdots x_r)=
 \sum_{\substack{\pi\text{ partition of }\{1,\ldots,r\}\\
                  |\pi|=n,\ D\ne\varnothing}}
       \prod_{D\in\pi}f\left(\prod_{j\in D}x_j\right).
\]
The expression is symmetric and multilinear, hence descends to
$\operatorname{Sym}^r(M)$ for every module $M$.
The $n$-fold convolution sums over ordered partitions; $f(1)=0$
removes empty blocks. Forgetting their order proves
$f^{*n}=n!\gamma_n(f)$ combinatorially, without division.
The sum axiom colors each block by either summand; the product
axiom chooses which of $m+n$ blocks belong to its first factor,
giving $\binom{m+n}{m}$; iteration groups $mn$ blocks into $m$
unordered groups of $n$, giving $c_{m,n}$. For the scalar axiom,
first partition the factors used by $f$ into $n$ nonempty blocks,
and distribute the remaining factors among the corresponding $n$
scalar factors. The two sides enumerate these same choices, with
the same products in $R$. These finite bijections verify all PD
axioms over $\mathbb Z$ and hence over every $k$.
\end{proof}

\paragraph{Finite verification and literature search.}
The exact program \nolinkurl{code/D_finite_pd_checks.py} and its
JSON output record all cardinalities, counts and elapsed time.
Finite index ranges supplement the all-index proofs above; they
do not replace them. Searches for convolution divided powers,
nonflat coalgebras, and square-zero Frobenius divided-power laws
found the opened flat comparison and the known square-zero
construction. No exact prior source for the second law in D011
was found in those searches; we claim no priority. The source and
search ledger records precisely what was opened.
